% =====================================================================
% Supportive-tool subsection template — STM32Cube Ecosystem Benchmark
% For a tool the AUTHOR BUILT THEMSELVES to enhance the benchmarking effort
% (the author's own engineering contribution — not a benchmark, and not a
% third-party tool merely used). Write it in the "we designed and built" voice.
% Copy into the methodology/tooling chapter, replace <...> placeholders,
% delete unused blocks. Keep labels meaningful and unique. pdfLaTeX-safe.
% Arc: problem -> design/architecture -> implementation -> integration ->
%      usage -> value/outcome.
% =====================================================================

\subsection{<Tool name, e.g. Automated flash-and-measure harness>}
\label{sec:<tool-key>}

% --- Problem: the benchmarking pain point this tool removes -----------
Running the benchmarks by hand <state the manual / error-prone / missing
capability> made <state the consequence: slow iteration, inconsistent
conditions, no traceability>. We therefore developed <tool> to <one-line
purpose>, so that <the benefit to the benchmarking effort>.

% --- Design & architecture: how it is structured ---------------------
\subsubsection{Design and architecture}
<Describe the components and how data flows from input to output.> The
overall structure is shown in Figure~\ref{fig:<tool-key>-arch}.

\begin{figure}[H]
  \centering
  % Replace with the diagram-designer's TikZ architecture/workflow figure,
  % or an image under Figures/<tool-key>/.
  \includegraphics[width=0.85\textwidth]{Figures/<tool-key>/<arch-file>}
  \caption{Architecture of <tool>: <input> $\rightarrow$ <processing> $\rightarrow$ <output>.}
  \label{fig:<tool-key>-arch}
\end{figure}

% --- Implementation: stack and key engineering choices ---------------
\subsubsection{Implementation}
We implemented <tool> in <language/framework>, relying on <key
libraries/external tools>~\cite{<lib-key>}. <Explain one or two notable
engineering decisions and why — e.g. how runs are timestamped, how the SWO/UART
stream is parsed, how configurations are swept.> The key parameters are
summarized in Table~\ref{tab:<tool-key>-io}; the full listing is given in
Appendix~\ref{chap:<appendix-label>} (Listing~\ref{lst:<tool-key>}).

% Inputs/outputs (or configuration) table
\begin{table}[H]
  \centering
  \begin{tabular}{l l}
    \toprule
    Aspect & Value \\
    \midrule
    Inputs        & <files / interfaces consumed> \\
    Outputs       & <artifacts produced: CSV/JSON/plots/report> \\
    Interfaces    & <UART / SWO / ITM / OpenOCD / CubeProgrammer> \\
    Stack         & <language, frameworks, key libraries> \\
    Invocation    & <CLI / GUI / CI step> \\
    \bottomrule
  \end{tabular}
  \caption{Inputs, outputs, and stack of <tool>.}
  \label{tab:<tool-key>-io}
\end{table}

% --- Integration: where it plugs into the benchmark workflow ---------
\subsubsection{Integration with the benchmarking workflow}
<Explain how the tool fits the pipeline (configure $\rightarrow$ build $\rightarrow$ flash
$\rightarrow$ measure $\rightarrow$ aggregate) and which benchmarks depend on it.>
The benchmarks of Chapter~\ref{chap:<results-chapter>} were produced with this tool,
ensuring identical conditions across configurations.

% --- Usage: how it is operated ---------------------------------------
\subsubsection{Usage}
<Show a representative invocation and what it produces. Keep it short; push the
full command reference / config to the appendix.>

% --- Value / outcome: the concrete improvement, honestly ------------
\subsubsection{Outcome}
<State what the tool improved for the benchmarking effort: reproducibility,
reduced per-run effort, consistent conditions, traceable logs. Give evidence
where available (e.g. a full sweep now runs unattended instead of manually).
Note limitations and what remains manual.>

\paragraph{Takeaway} <One or two sentences: the engineering value the tool adds
to the report's benchmarking methodology.>
